\section{Prompt Templates}
\label{sec:appendix_prompts}

We provide representative prompt templates for the \gsm benchmark. Prompts for other benchmarks follow the same structure with domain-appropriate exemplars.

\subsection{\fewshotcot Exemplar (GSM8K)}

\begin{lstlisting}
Q: There are 15 trees in the grove. Grove workers will plant
trees in the grove today. After they are done, there will be 21
trees. How many trees did the grove workers plant today?
A: Step 1: We start with 15 trees.
Step 2: After planting, there are 21 trees.
Step 3: The number planted is 21 - 15 = 6.
The answer is 6.
\end{lstlisting}

\subsection{\storycot Exemplar (GSM8K)}

\begin{lstlisting}
Q: There are 15 trees in the grove. Grove workers will plant
trees in the grove today. After they are done, there will be 21
trees. How many trees did the grove workers plant today?
A: Imagine a grove at dawn with 15 trees standing tall. A team
of workers arrives with saplings on a truck. They spend the
morning digging holes and carefully planting new trees among the
existing ones. By sunset, someone walks through and counts every
tree -- there are now 21 in total. The original 15 haven't
changed, so the workers must have added 21 minus 15, which is 6
new trees during their day's work.
The answer is 6.
\end{lstlisting}

\subsection{\storycot Exemplar (CommonsenseQA)}

\begin{lstlisting}
Q: Where do you put your grapes just before checking out?
(A) mouth (B) move (C) basket (D) fruit stand (E) fruit market
A: Picture yourself in a grocery store, pushing a cart through
the produce aisle. You spot a display of fresh grapes and pick
up a bunch, examining them briefly. Before heading to the
checkout counter, you place the grapes into your shopping
basket alongside your other items.
The answer is C.
\end{lstlisting}

\subsection{\storycot Exemplar (ARC-Challenge)}

\begin{lstlisting}
Q: Which property of a mineral can be determined just by
looking at it?
(A) luster (B) mass (C) weight (D) hardness
A: Imagine a geologist crouching beside a rocky outcrop on a
sunny afternoon. She picks up a mineral specimen and holds it up
to the light, tilting it back and forth. The way the surface
catches and reflects the sunlight -- its sheen and brilliance --
tells her immediately about its luster. She would need a scale
for mass, a balance for weight, and a scratch test for hardness,
but luster reveals itself to the naked eye.
The answer is A.
\end{lstlisting}
